\documentclass[10pt,a4paper]{article}
\usepackage{amsmath,amssymb,amsthm,graphicx}
\usepackage{amsmath,amssymb,amsthm,hyperref,xfrac}
\providecommand{\ie}{i.e.\ }
\providecommand{\skpt}{\smallskip}
\usepackage[scr=rsfs,cal=euler]{mathalfa}
\multlinegap0pt

\hyphenation{begin decom-po-si-tion sequence}
\newenvironment{enumeratea}
{\bgroup\def\theenumi{\alph{enumi}}\begin{enumerate}}
{\end{enumerate}\egroup}

\newcommand{\ba}{{\boldsymbol a}}
\newcommand{\bb}{{\boldsymbol b}}
\newcommand{\bd}{{\boldsymbol d}}
\newcommand{\bu}{{\boldsymbol u}}
\newcommand{\bx}{{\boldsymbol x}}
\newcommand{\bmy}{{\boldsymbol y}}

\newcommand{\bD}{{\boldsymbol D}}
\newcommand{\bL}{{\boldsymbol L}}
\newcommand{\bP}{{\boldsymbol P}}
\newcommand{\bV}{{\boldsymbol V}}
\newcommand{\bX}{{\boldsymbol X}}
\newcommand{\bZ}{{\boldsymbol Z}}

\newcommand{\Psfrac}[2]{(\sfrac{#1}{#2})}
\newcommand{\psfrac}[2]{\sfrac{(#1)}{#2}}
\newcommand{\bigO}[1]{O(#1)}

\usepackage{mathtools}
\usepackage{bbm}

\newcommand{\NN}{\mathbb{N}}
\newcommand{\ZZ}{\mathbb{Z}}
\newcommand{\RR}{\mathbb{R}}
\newcommand{\CC}{\mathbb{C}}
\newcommand{\PP}{\mathbb{P}}
\newcommand{\EE}{\mathbb{E}}

\newcommand{\boldalpha}{\boldsymbol{\alpha}}
\newcommand{\bolddelta}{\boldsymbol{\delta}}
\newcommand{\boldrho}{\boldsymbol{\rho}}
\newcommand{\dee}{\,\mathrm{d}}
\newcommand{\one}{\mathbbm{1}}
\newcommand{\BC}{\mathcal{B}}
\newcommand{\EC}{\mathcal{E}}
\newcommand{\RS}{\mathscr{R}}
\newcommand{\prob}[1]{\PP[ #1 ]}
\newcommand{\bigoo}[1]{O\!\left(#1 \right)}
\newcommand{\abs}[1]{\left\lvert #1 \right\rvert}
\newcommand{\floor}[1]{\left\lfloor #1 \right\rfloor}
\newcommand{\riemint}[4]{\int_{#1}^{#2}\! #3 \, \mathrm{d}#4}
\newcommand{\maxnorm}[1]{\left\| #1 \right\|_{\infty}}

\renewcommand{\Re}{\textup{Re}}
\renewcommand{\Im}{\textup{Im}}

\newcommand{\dtv}{\mathrm{d}_{\mathrm{TV}}}
\newcommand{\dw}{\mathrm{d}_{\mathrm W}}

\let\epsilon\varepsilon

\newcommand{\pextra}{P_{\mathtt{extra}}}
\newcommand{\parratia}{P_{\mathtt{Arratia}}}

\newcommand{\simplex}{\Delta^{k-1}}

\newcommand{\bg}{\big}
\newcommand{\bgg}{\Big}
\newcommand{\bggg}{\bigg}

\newcommand{\Rnt}{R_{\mathrm{NT}}}
\newcommand{\Rprob}{R_{\mathrm{PD}}}
\newcommand{\ECnt}{\EC_{\mathrm{NT}}}
\newcommand{\ECprob}{\EC_{\mathrm{PD}}}

\newcommand{\Bnt}{\mathcal{B}_{\mathrm{NT}}}
\newcommand{\Bprob}{\mathcal{B}_{\mathrm{PD}}}

\DeclareMathOperator{\vard}{Var}
\DeclareMathOperator{\Dir}{Dir}

\theoremstyle{plain}
\newtheorem{cor}{Corollary}[section]
\newtheorem{lem}[cor]{Lemma}
\newtheorem{prop}[cor]{Proposition}
\newtheorem{conj}[cor]{Conjecture}
\newtheorem{theo}{Theorem}

\theoremstyle{definition}
\newtheorem{example}{Example}
\newtheorem{deff}[cor]{Definition}
\newtheorem{rem}{Remark}[section]
\newtheorem*{rems}{Remark}
\newtheorem*{remss}{Remarks}

\makeatletter
\ams@newcommand{\multiint}[1]{\DOTSI\protect\MultiIntegral{#1}}
\renewcommand{\MultiIntegral}[1]{%
\edef\ints@c{\noexpand\intop
\ifnum#1=\z@\noexpand\intdots@\else\noexpand\intkern@\fi
\replicate{#1-2}{\noexpand\intop\noexpand\intkern@}%
\noexpand\intop
\noexpand\ilimits@
}%
\futurelet\@let@token\ints@a
}
\makeatother
\ExplSyntaxOn
\cs_new:Npn \replicate #1 #2 { \prg_replicate:nn { #1 } { #2 } }
\ExplSyntaxOff
\title{The probability that the second or the third event fails is of order one over the logarithm}
\author{Tony Haddad, Dimitris Koukoulopoulos}
\date{doi:10.5802/jep.317; CC BY 4.0}
\begin{document}
\maketitle
\noindent\textit{Editorial note.} The counted statement is the author's labelled lem environment at source0.tex lines 963-966; its proof environment opens at line 968, the line after the blank separator, so the pairing is by adjacency and the retained proof range 968-1041 is the whole environment.\par\medskip
\noindent\textit{Editorial note.} Three environments are retained as context because the counted proof depends on them: the exponential law of the truncated weighted sum prop:icexponential (source lines 814-817), applied at source line 984; the decomposition lemma lem:decomposition J neq J* (source lines 842-848), used at source line 968; and the numbered auxiliary expectation eq:E(z) (source lines 923-926), referred to at source line 973.\par\medskip
\noindent\textit{Editorial note.} This article carries no label-number dump and the retained counter values are left empty for the whole source, so the printed numbering of the card is the wrapper's default; the identity of the counted item is its author label lem:prob e2c e3c.\par\medskip
\noindent\textit{Editorial scope.} One original result of the pinned publisher article, counted once: the estimate of order one over the logarithm for the failure of the second and third coupling events (author label lem:prob e2c e3c).\par\medskip
\noindent\textit{Allowed standard prerequisites (editorial).}\begin{itemize}
\item The three events of the coupling. The events E1, E2 and E3 control respectively the total mass, the absence of a large ratio between consecutive points and the separation of the point process from its shifted copy.. They are the events whose complements are estimated in the statement.
\item Mecke equation. For a Poisson process the expectation of a sum over the points of a function of the point and of the process with that point removed is an integral of that function against the intensity measure.. It turns the expectation of the sum over pairs of points into a multiple integral.
\item Exponential law of the truncated weighted sum. The sum of the second coordinates of the points whose first coordinate exceeds a level follows an exponential law with parameter equal to that level.. It supplies the tail bound on the probability appearing in the integral.
\end{itemize}\medskip
\begin{prop}[Arratia, \cite{Arratia02}]
\label{prop:icexponential}
For any fixed $t > 0$, the random variable $I_t$ follows an exponential distribution of parameter $t$, \ie $\prob{I_t > y} = e^{-ty}$ for all $y > 0$.
\end{prop}

\begin{lem}
\label{lem:decomposition J neq J*}
For $x > 1$, we~have $\EC_1\cap \EC_2\cap \EC_3\subseteq \{ J_x = J_x^*\}$. In~particular,
\[
\PP\bg[J_x \ne J_x^*\bg] \le \PP[\EC_1^c] + \PP[\EC_1 \cap \EC_2^c] + \PP[\EC_1 \cap \EC_3^c].
\]
\end{lem}

\begin{equation}
\label{eq:E(z)}
E(z) \coloneqq \EE \bggg[\mathop{\sum\sum\sum}_{\substack{i, j, k \in \ZZ:\ i>j>k}}r(Y_i)r(Y_j) \cdot \one_{|S_k-\log x|\le z} \cdot \one_{W_k \le (\log_2x)^{-10}} \bgg] \ll \frac{z}{\log x}
\end{equation}

\begin{lem}
\label{lem:prob e2c e3c}
We have $\PP[\EC_1 \cap \EC_2^c] + \PP[\EC_1 \cap \EC_3^c] \ll 1/\log x$ for $x\ge2$.
\end{lem}

\begin{proof} Recall that $\EC_2$ failing means that there exists $i\ge1$ such that $W_i/W_0\le 1+\eta/\min\{Y_0, Y_i\}^2$. In~addition, recall that the event $\EC_1$ implies that $Y_0 > e^{-\gamma}$ (this was explained in the beginning of the proof of Lemma \ref{lem:decomposition J neq J*}). Hence, we~have
\[
\PP[\EC_1 \cap \EC_2^c] \le
\EE\bggg[ \mathop{\sum\sum}_{\substack{(W,Y) \in \RS \\ (W',Y')\in\RS}} \one_{1< \frac{W'}{W} \le 1+\sfrac{\eta}{\min\{Y,Y'\}^2}}\cdot\one_{Y',Y>e^{-\gamma}} \cdot \one_{I_{W, W'} + Y'\in(\log x-Y,\log x]} \bggg]
\]
with $I_{t, t'} \coloneqq \sum_{(W, Y) \in \RS,\ W \in \RR_{>t}\setminus\{t'\}} Y$. We~use Mecke's equation as in the proof of \eqref{eq:E(z)} to get that the expectation above is equal to
\[
\iiiint\limits_{\substack{0 < w < w' < w(1+\eta/\min\{y,y'\}^2)\\ y, y' > e^{-\gamma},\ y'\le \log x}} \PP\bg[\log x - y<I_w + y' \le \log x\bg]\cdot e^{-wy-w'y'}
\dee w\dee w'\dee y\dee y'.
\]
We have $e^{-w'y'}\le e^{-wy'}$, thus
\[
\PP[\EC_1 \cap \EC_2^c] \ll
\iiint\limits_{\substack{w>0,\ y, y' > e^{-\gamma} \\ y'\le \log x}}
\PP\bg[\log x - y<I_w + y' \le \log x\bg]\cdot \frac{we^{-w(y+y')}}{\min\{y,y'\}^2} \dee w\dee y\dee y'.
\]
By Proposition \ref{prop:icexponential}, we~have
\[
\PP\bg[\log x - y<I_w + y' \le \log x\bg] \le \begin{cases}
e^{-w(\log x - y-y')}&\text{if $y \le \log x$,} \\
1 & \text{if $y > \log x$.}
\end{cases}
\]
Therefore,
\begin{align*}
\PP[\EC_1 \cap \EC_2^c]
&\ll \iiint\limits_{\substack{w>0,\ y, y' > e^{-\gamma} \\ y,y'\le \log x}}
\frac{we^{-w\log x}}{\min\{y,y'\}^2}\dee w\dee y\dee y'
+\iiint\limits_{\substack{w>0,\ y, y' > e^{-\gamma} \\ y'\le \log x<y}}
\frac{we^{-w(y+y')}}{(y')^2}\dee w\dee y\dee y' \\
&\le \iint\limits_{e^{-\gamma}<y, y'\le \log x}
\frac{1}{\min\{y,y'\}^2(\log x)^2}\dee y\dee y'
+\iint\limits_{\substack{y, y' > e^{-\gamma} \\ y'\le \log x<y}}
\frac{1}{(yy')^2} \dee y\dee y' \\
&\ll 1/\log x.
\end{align*}
This completes the proof of the claimed bound on $\PP[\EC_1\cap \EC_2^c]$.

Finally, we~bound $\PP[\EC_1 \cap \EC_3^c]$ using a very similar argument. We~have
\[
\begin{split}
&\PP[\EC_1 \cap \EC_3^c] \\
&\qquad \le
\EE\bggg[ \mathop{\sum\sum}_{(W,Y),(W',Y') \in\RS} \one_{1<\frac{W}{W'}\le 1+\frac{\eta}{\min\{Y,Y'\}^2}}\cdot\one_{Y',Y>e^{-\gamma}} \cdot \one_{I_W\in(\log x-Y,\log x]} \bggg] \\
&\qquad =
\iiiint\limits_{\substack{0 < w'<w < w'(1+\eta/\min\{y,y'\}^2)\\ y, y' > e^{-\gamma}}} \PP\bg[\log x - y<I_w \le \log x\bg]\cdot e^{-wy-w'y'}
\dee w\dee w'\dee y\dee y' \\
&\qquad \le J_1+J_2,
\end{split}
\]
where
\[
J_1 \coloneqq \iiiint\limits_{\substack{0 < w'<w < w'(1+\eta/\min\{y,y'\}^2)\\ y, y' > e^{-\gamma},\ y\le \log x}} e^{-w\log x-w'y'}
\dee w\dee w'\dee y\dee y'
\]
and
\[
J_2 \coloneqq \iiiint\limits_{\substack{0 < w'<w < w'(1+\eta/\min\{y,y'\}^2)\\ y>\log x,\ y' > e^{-\gamma}}} e^{-wy-w'y'}
\dee w\dee w'\dee y\dee y'.
\]
Using $e^{-w\log x}\le e^{-w'\log x}$, we~find
\begin{align*}
J_1&\ll \iiint\limits_{\substack{w'>0,\ y, y' > e^{-\gamma}\\ y\le \log x}} \frac{w'e^{-w'(y'+\log x)}}{\min\{y,y'\}^2} \dee w'\dee y\dee y' \\
&= \iint\limits_{\substack{y, y' > e^{-\gamma} \\ y\le \log x}} \frac{1}{\min\{y,y'\}^2(y'+\log x)^2} \dee y\dee y'\\
& \ll1/\log x.
\end{align*}
Similarly, we~have
\begin{align*}
J_2&\ll \iiint\limits_{\substack{w'>0,\ y' > e^{-\gamma}\\ y>\log x}} \frac{w'e^{-w'(y+y')} }{\min\{y,y'\}^2} \dee w'\dee y\dee y' \\
&= \iint\limits_{\substack{y' > e^{-\gamma} \\ y>\log x}} \frac{1}{\min\{y,y'\}^2(y+y')^2} \dee y\dee y'\\
& \ll1/\log x.
\end{align*}
This implies that $\PP[\EC_1\cap\EC_3^c]\ll 1/\log x$, thus completing the proof of the lemma.
\end{proof}

\begin{thebibliography}{99}
\bibitem{Arratia98} R. Arratia - “On the central role of scale invariant Poisson processes on (0,∞)”, in Microsurveys in discrete probability (Princeton, NJ, 1997), DIMACS Ser. Discrete Math. Theoret. Comput. Sci., vol. 41, American Mathematical Society, Providence, RI, 1998, p. 21-41
\bibitem{Arratia02} R. Arratia - “On the amount of dependence in the prime factorization of a uniform random integer”, in Contemporary combinatorics, Bolyai Soc. Math. Stud., vol. 10, János Bolyai Math. Soc., Budapest, 2002, p. 29-91
\bibitem{ABT06} R. Arratia, A. D. Barbour \& S. Tavaré - “A tale of three couplings: Poisson-Dirichlet and GEM approximations for random permutations”, Combin. Probab. Comput. 15 (2006) no. 1-2, p. 31-62
\bibitem{BareikisManstavicius07} G. Bareikis \& E. Manstavičius - “On the DDT theorem”, Acta Arith. 126 (2007) no. 2, p. 155-168
\bibitem{Bill72} P. Billingsley - “On the distribution of large prime divisors”, Period. Math. Hungar. 2 (1972), p. 283-289
\bibitem{delaBretecheTenenbaum16} R. de la Bretèche \& G. Tenenbaum - “Sur les processus arithmétiques liés aux diviseurs”, Adv. in Appl. Probab. 48 (2016) no. A, p. 63-76
\bibitem{DDT79} J.-M. Deshouillers, F. Dress \& G. Tenenbaum - “Lois de répartition des diviseurs. I”, Acta Arith. 34 (1979) no. 4, p. 273-285 (loose errata)
\bibitem{DonnellyTavare87} P. Donnelly \& S. Tavaré - “The population genealogy of the infinitely-many neutral alleles model”, J. Math. Biol. 25 (1987) no. 4, p. 381-391
\bibitem{Feng10} S. Feng - The Poisson-Dirichlet distribution and related topics. Models and asymptotic behaviors, Probability and its Appl., Springer, Heidelberg, 2010
\bibitem{JKB} N. L. Johnson, S. Kotz \& N. Balakrishnan - Discrete multivariate distributions, Wiley Series in Probab. and Stat.: Applied Probab. and Stat., John Wiley \& Sons, Inc., New York, 1997
\bibitem{Kingman93} J. F. C. Kingman - Poisson processes, Oxford Studies in Probab., vol. 3, The Clarendon Press, Oxford University Press, New York, 1993
\bibitem{KotzBal00} S. Kotz, N. Balakrishnan \& N. L. Johnson - Continuous multivariate distributions. Vol. 1. Models and applications, Wiley Series in Probab. and Stat.: Applied Probab. and Stat., Wiley-Interscience, New York, 2000
\bibitem{Koukoulo19} D. Koukoulopoulos - The distribution of prime numbers, Graduate Studies in Math., vol. 203, American Mathematical Society, Providence, RI, 2019
\bibitem{LastPenrose18} G. Last \& M. Penrose - Lectures on the Poisson process, Institute of Math. Stat. Textbooks, vol. 7, Cambridge University Press, Cambridge, 2018
\bibitem{Ken22} S.-K. Leung - “Dirichlet law for factorisation of integers, polynomials and permutations”, Math. Proc. Cambridge Philos. Soc. 175 (2023) no. 3, p. 649-676
\bibitem{KenWeb24} S.-K. Leung - “Dirichlet law with arbitrary parameters for factorizations of integers” (2025), Retrieved on 24 August 2025 from https://sites.google.com/view/sunkaileung/miscellaneous
\bibitem{MV07} H. L. Montgomery \& R. C. Vaughan - Multiplicative number theory. I. Classical theory, Cambridge Studies in Advanced Math., vol. 97, Cambridge University Press, Cambridge, 2007
\bibitem{NyandwiSmati13} S. Nyandwi \& A. Smati - “Distribution laws of pairs of divisors”, Integers 13 (2013), article ID A13, 13 pages
\bibitem{Tenenbaum00} G. Tenenbaum - “A rate estimate in Billingsley’s theorem for the size distribution of large prime factors”, Q. J. Math. 51 (2000) no. 3, p. 385-403
\end{thebibliography}
\end{document}
